% GENERATED FILE. Edit canonical lessons, then run npm run generate:books.
% canonical-source-hash: fnv1a64:e57354644fede142
% canonical-lessons: ES-C06-estudiar

\chapter{Estudiar --- Studying Is Being Eager}
\label{ch:spanish-estudiar}
\hlchaptermodality{21}

\begin{chapteropening}
\textbf{By the end of this chapter:} \emph{I can say what I study, with a third regular -ar verb and no new pattern.}

studere, to be eager - the root under student, studio and studious.
\end{chapteropening}

\section[estudiar]{estudiar — \textquotedblleft{}to study,\textquotedblright{} and the eagerness inside it}

\label{lesson:ES-C06-estudiar}



\begin{quote}
Your third \textbf{-ar} verb. By now the pattern should feel automatic — that's exactly what we want. \emph{estudiar} also hides a nicer feeling than the classroom suggests: at root it means \textbf{eagerness.}
\end{quote}

\begin{sounds}
\begin{itemize}
\raggedright
  \item \texttt{st-no-e} — Spanish props up \emph{st-} with an \emph{e}: \textbf{es-tudiar} (the same reflex behind \emph{estar}, \emph{España}, \emph{escuela}). \emph{es-too-DYAR}.
  \item \texttt{diphthong-io} — the \emph{-iar} glides: \emph{es-too-\textbf{DYAR}}.
\end{itemize}
\end{sounds}

\begin{cousinweb}
\textbf{estudiar} comes from Latin \textbf{studēre}, \emph{\textquotedblleft{}to be eager, to apply oneself zealously.\textquotedblright{}} A \emph{studium} was \textbf{zeal, enthusiasm, devotion} — only later \textquotedblleft{}study.\textquotedblright{} So to \emph{study} is, at heart, to \emph{throw yourself eagerly} at something.

That \textbf{studēre} is everywhere in English:

\begin{itemize}
\raggedright
  \item \textbf{student} — literally \textquotedblleft{}one who is \emph{eager}.\textquotedblright{}
  \item \textbf{studio}, \textbf{studious}, \textbf{study} — the place, the trait, the act.
  \item Italian \textbf{studio} and the musical \emph{studio} keep the \textquotedblleft{}devoted practice\textquotedblright{} sense.
\end{itemize}

So \emph{Estudio español} is \textquotedblleft{}I \emph{devote myself eagerly} to Spanish\textquotedblright{} — a nicer picture than mere schoolwork.
\end{cousinweb}

\begin{grammarlens}[title={Grammar Lens: the template holds}]
Drop \textbf{-ar}, add the ending — for the third time:

\noindent
\begin{tabularx}{\linewidth}{@{}>{\raggedright\arraybackslash}X>{\raggedright\arraybackslash}X>{\raggedright\arraybackslash}X@{}}
\toprule
\textbf{person} & \textbf{form} & \textbf{meaning} \\
\midrule
I (ending carries the subject) & \textbf{estudio} & I study \\
tú & \textbf{estudias} & you study \\
usted & \textbf{estudia} & you (formal) study \\
\bottomrule
\end{tabularx}

Three verbs, one pattern: \emph{hablo, trabajo, estudio}. The singular regular \textbf{-ar} present is now reusable rather than memorized verb by verb.
\end{grammarlens}

\subsection*{Your turn}
Say these aloud:

\begin{itemize}
\raggedright
  \item \textquotedblleft{}estudiar\textquotedblright{} — \emph{es-too-DYAR}
  \item \textquotedblleft{}estudio, estudias, estudia\textquotedblright{} — the now-familiar endings
  \item \textquotedblleft{}estudiar\textquotedblright{} then English \textquotedblleft{}student, studio\textquotedblright{} — the \emph{studēre} \textquotedblleft{}eagerness\textquotedblright{}
\end{itemize}

\begin{tcolorbox}[breakable,colback=teal!4,colframe=teal!35!black,title={Before you move on}]
What did \emph{studēre} originally mean? (\textquotedblleft{}To be eager / apply oneself.\textquotedblright{}) Its English cousins? (Student, studio, studious.) Say \textquotedblleft{}I study\textquotedblright{} — and notice it uses the same \emph{-o} as \emph{hablo} and \emph{trabajo}. (\emph{Estudio}.) Next: put a verb and a noun together into your \textbf{first real sentence}.
\end{tcolorbox}
